% ============================================================================
%  spiral_v16_preprint.tex   --  methodology report, v16 line
%
%  Auditing the auditor: what a verification ledger found when it was turned
%  on its own author.
%
%  Framing (must hold): this is a NUMERICAL-EXPERIMENT / METHODOLOGY report.
%  It claims NO physical result. Every number is transcribed from a frozen
%  reference run or from a registered, dated log entry; each is annotated with
%  its source. Numbers that this report could not verify are NOT stated.
%
%  Reference run:  v16/_runall_v16.3_r1.log  (EXIT=0, 1119.7 s, 8 cores)
%  Companion:      v15/spiral_v15_preprint.tex (the ledger's first paper)
%
%  Self-contained: references are inline (\thebibliography), so a single
%  `pdflatex spiral_v16_preprint.tex` pass suffices; no bibtex run needed.
% ============================================================================
\documentclass[11pt,a4paper]{article}

\usepackage[utf8]{inputenc}
\usepackage[T1]{fontenc}
\usepackage{amsmath,amssymb}
\usepackage{booktabs}
\usepackage{array}
\usepackage{underscore}
\usepackage[margin=2.4cm]{geometry}
\usepackage[hidelinks]{hyperref}
\usepackage{enumitem}
\usepackage{xcolor}

\newcommand{\code}[1]{\texttt{\small #1}}
\newcommand{\src}[1]{{\footnotesize\texttt{#1}}}

\title{\bfseries Auditing the auditor \\[2pt]
\large --- what a verification ledger found when it was \\
\large turned on its own author --- \\[8pt]
\normalsize demonstrated on an Ising / MERA / reaction--diffusion pipeline}

\author{%
   JianZhong Yan\\
  \small \textit{(independent researcher)}\\
  \small \texttt{mohedanovinita984@gmail.com}\\
}
\date{\today}

\begin{document}
\maketitle

% ---------------------------------------------------------------------------
\begin{abstract}
\noindent
A companion paper introduced an operational criterion for a genuine
\emph{negative control} in multi-stage numerical pipelines --- \textbf{it must be
able to fail when the conclusion does not hold} --- and built a per-layer ledger
from it. This paper reports what happened when the same discipline was turned on
the ledger's own author and code.

\textbf{We report five self-audit outcomes.} (i) The ledger's last coverage gap
was \emph{closed} (negative-control coverage $6/7 \to 7/7$ layers) and the layer
was \emph{still not promoted} --- because the new control shares its construction
with the criterion it guards, so it calibrates a sensitivity scale rather than
supplying independent falsification. (ii) A registered ``derivation'' of a
uniqueness statement was found to be a \emph{tautology}: two ready-made theorems
juxtaposed, with the entropy object in the claim not the entropy the code
computes. (iii) Two registrations were found to be \emph{stale} --- the code did
not do what the registry said --- one of which concerned the pass-rate
\emph{denominator} itself. (iv) A refuted inference was found to have been
written in \emph{four} places; it was swept repo-wide rather than patched where
it surfaced. (v) A guard too weak to fail printed a
\emph{plausible-looking fake reading} --- the very failure mode the method
exists to prevent, occurring live inside the method's own reference
implementation.

Result (v) is the paper's centrepiece: with the imaginary-time operator
$e^{-\tau H}$, the norm grows as $e^{\tau|E_0|}$; at $\tau{=}20$ the vector's
entries are still finite ($\sim\!10^{176}$) while \code{np.linalg.norm}
overflows \emph{inside its own squaring}, so normalising by it silently yields
the zero vector and the code prints $S(L/2)=0$ with overlap $0$ --- a reading
indistinguishable, in form, from a real one. The fix is to guard on the maximum
entry rather than the norm.

\textbf{Ledger state.} 27 metrics, 78 guards, 13 negative controls; scored guard
pass rate $57/57=100.0\%$; 21 checks registered as \emph{diagnostics} and
excluded from the denominator, 16 of which do not pass by design.
\textbf{No layer was promoted}: the maturity vector is unchanged from the
companion paper.

\smallskip
\noindent\textbf{Boundary.} ``$100\%$'' is a statement about \emph{guard
execution}, not about the truth of any physical emergence claim. This paper
claims no physical result beyond textbook Ising criticality, and its one
apparently positive template (Sec.~\ref{sec:positive}) is a statement about a
\emph{criterion}, not about a spacetime.

\smallskip
\noindent\textbf{Code availability:}
\url{https://github.com/yanjianzhong/spiral-emergence}
\end{abstract}

% ---------------------------------------------------------------------------
\section{Introduction}
\label{sec:intro}

\subsection{Recapitulation in one paragraph}

The companion paper \cite{v15paper} observed that the assertions enforcing
correctness at each stage of a layered numerical pipeline are, in practice,
\emph{positivity}, \emph{finiteness} and \emph{non-emptiness} checks --- and that
such checks are \emph{consistent with any conclusion} \cite{popper}. ``The
Schmidt rank is $\ge 2$'' holds for a correct calculation and for a calculation
wrong in every way that matters. It proposed instead an operational criterion: a
genuine negative control \textbf{must be able to fail when the conclusion does
not hold}, and a ledger built from controls satisfying it, carrying a
\emph{completeness guard} that checks the ledger's own citations against the
\emph{runtime} registry rather than trusting the ledger's self-report.

That paper reported a ledger covering 6 of 7 layers, and three negative
results.

\subsection{What is new here}

This paper is not a re-derivation. It reports the outcome of the next
iteration, in which the discipline was aimed at the artefact doing the
auditing: the registry, the ledger, the documentation, and the author's own
recorded inferences.

The result is not a new physical claim. It is a catalogue of the ways a
verification discipline \emph{actually} fails in practice, each instance
observed rather than hypothesised --- and, most usefully, an instance in which
the failure was caught by the discipline itself rather than by an external
reviewer. We take that last class to be the strongest available evidence that a
verification method is doing work: a method that has never caught its own
author is indistinguishable from a method that cannot.

\subsection{Contributions}

\begin{enumerate}[leftmargin=1.4em,itemsep=2pt]
  \item \textbf{A runtime partition} of one check registry into scored checks
        and diagnostics (Sec.~\ref{sec:split}), replacing the fragile pattern
        of two hand-maintained lists. The partition is computed at report time
        from a single list, so the denominator cannot drift from the registry.
  \item \textbf{Five documented self-audit outcomes} (Sec.~\ref{sec:selfaudit}),
        each with the rule it generalises to.
  \item \textbf{A live counter-example} to the method's own thesis
        (Sec.~\ref{sec:fake}): a guard that passed while the value it guarded
        was meaningless, producing a reading of exactly the plausible kind the
        method is designed to expose.
  \item \textbf{An honest boundary on aggregate pass rates}
        (Sec.~\ref{sec:aggregate}): on this pipeline a single ``metric pass
        rate'' is not well defined, because one tier carries no physical
        target line at all. Quoting one would be a category error.
  \item \textbf{One positive template} (Sec.~\ref{sec:positive}) --- the only
        place in the project where a criterion \emph{selected} a convention
        rather than being determined by it, offered as the constructive
        counterpart to the five negative results.
\end{enumerate}

% ---------------------------------------------------------------------------
\section{Substrate and the frozen reference run}
\label{sec:substrate}

\subsection{The pipeline}

The substrate is a seven-layer pipeline ($L1 \dots L7$ plus an external anchor),
unchanged from the companion paper: a void/chain construction, a critical Ising
break, MERA entanglement and central-charge extraction \cite{vidal,calabrese}, a
graph-geometric readout on a boundary-correlation graph \cite{ollivier,forman},
a cMERA branch \cite{haegeman}, a Gray--Scott reaction--diffusion system
\cite{pearson}, and an external data anchor. The exact diagonalisation, the
Jordan--Wigner closed form for the periodic transverse-field Ising chain
\cite{pfeuty}, and the MERA and cMERA machinery are as described there; this
paper alters no earlier readout, threshold, criterion or denominator.

\subsection{The two ledgers must not be mixed}

The single most error-prone feature of the reporting is that the pipeline keeps
\emph{two parallel ledgers}:

\begin{itemize}[leftmargin=1.4em,itemsep=2pt]
  \item \textbf{Metrics} --- registered by \code{record_metric(..., passed=)}:
        physical verdicts carrying a threshold (``does this number meet its
        target?'').
  \item \textbf{Guards} --- registered by \code{record_guard(..., ok=)}:
        falsifiable assertions (``does this statement hold?'').
\end{itemize}

The two are disjoint and must not be summed. A single composite ``score'' across
both would be meaningless, and a report that presented one would be doing the
thing this paper warns about.

\subsection{Frozen reference run}

All v16 numbers below come from one frozen run of the repository's single
runner, archived as \src{v16/\_runall\_v16.3\_r1.log}.

\begin{table}[h]
\centering
\small
\begin{tabular}{l l l}
\toprule
 & companion paper & \textbf{this paper} \\
\midrule
exit code                    & \code{EXIT=0}      & \code{EXIT=0} \\
scored guards                & $54/54 = 100.0\%$  & $\mathbf{57/57 = 100.0\%}$ \\
metrics (registry size)      & 18                 & \textbf{27} \\
diagnostics (excluded)       & 20 (17 not passing)& \textbf{21 (16 not passing)} \\
negative-control coverage    & $6/7$ layers, gap \code{['L1']}
                                                    & $\mathbf{7/7}$ \textbf{layers, gap} \code{[]} \\
negative controls            & 12                 & \textbf{13} \\
wall time                    & 1245\,s            & \textbf{1119.7\,s} \\
\bottomrule
\end{tabular}
\caption{The ledger across two iterations. Row~3 gives \emph{registry size}, not
a pass rate --- see Sec.~\ref{sec:aggregate} for why the latter is not well
defined here. Guard counts are verbatim from \src{\_runall\_v16.3\_r1.log:1230}
and the registered sizes from \src{:1239}; wall time from \src{:1266}.}
\label{tab:ledger}
\end{table}

\subsection{Why an aggregate metric pass rate is not well defined}
\label{sec:aggregate}

It is tempting to write ``$27/27$ metrics at target'' by analogy with the
scored-guard rate. We do not, for a structural reason recorded in the audit
table: the metric registry spans three tiers, and the \textbf{third tier
carries no physical target line} --- its registered distribution is ``$0/0$
met'', with six further items held as diagnostics. A ratio whose denominator is
empty in one tier has no unique aggregate interpretation, and choosing one after
the fact would be exactly the kind of convention-driven reporting this project
exists to prevent. The per-tier report is emitted as a 20-panel figure and a
JSON artifact instead.

We record this as a \emph{methodological} result, not a shortcoming: a
verification report should state which of its aggregates are well defined, and
decline the rest.

% ---------------------------------------------------------------------------
\section{A runtime partition instead of two lists}
\label{sec:split}

\subsection{Mechanism}

Earlier iterations maintained scored checks and diagnostics as two separate
lists, with the pass-rate denominator hard-coded from one of them. This is
fragile: moving a check between the lists requires editing both the list and
the denominator, and a mismatch is silent --- the report still prints a
percentage, just a wrong one.

v16 replaces this with a single registry and a partition computed at report
time:

\begin{quote}\small
\code{guard_pass_rate()} iterates the \emph{same} \code{GUARDS} list and splits
it by the per-item flag \code{expect_pass}: items with \code{expect_pass=True}
form the scored set; the remainder are diagnostics.
\hfill\src{v16/\_metric/registry.py:153--161}
\end{quote}

The denominator is therefore not a number anyone can transcribe incorrectly: it
is derived from the same object the checks were recorded into.

\subsection{Why this is a control and not a convenience}

The partition is what makes the third column of Table~\ref{tab:ledger}
\emph{auditable}. In the frozen run the report line is verbatim:

\begin{quote}\small
\code{The pass rate for guards is 57/57 = 100.00\% (excluding 21 diagnostic items, of which 16 failed)}
\hfill\src{v16/\_runall\_v16.3\_r1.log:1230}
\end{quote}

so $57 + 21 = 78$ is not a claim about the code, it is a consequence of it. This
distinction turned out to matter immediately, and is the subject of
Sec.~\ref{sec:stale}.

% ---------------------------------------------------------------------------
\section{Five self-audit outcomes}
\label{sec:selfaudit}

Table~\ref{tab:selfaudit} summarises. Each subsection states the observation and
the rule it generalises to; Sec.~\ref{sec:commontheme} extracts what the five
share.

\begin{table}[h]
\centering
\small
\begin{tabular}{p{1.9cm} p{6.4cm} p{5.5cm}}
\toprule
 & \textbf{what was caught} & \textbf{disposition} \\
\midrule
\S\ref{sec:gapclosure} gap closure
 & coverage went $6/7 \to 7/7$, but the new control shares construction with
   the criterion it guards
 & layer \textbf{not} promoted; the control calibrates \emph{sensitivity},
   not independent falsification \\
\addlinespace
\S\ref{sec:tautology} tautology
 & a registered ``derivation'' was two juxtaposed theorems; the entropy object
   in the claim was not the entropy the code computes
 & registered as a \emph{tautology}; the gap it was meant to close is
   \textbf{not} narrowed \\
\addlinespace
\S\ref{sec:stale} stale registration
 & two registry entries described behaviour the code does not have; one touched
   the pass-rate denominator itself
 & corrected in place, original retained and annotated; denominator settled
   from code, not from a log line \\
\addlinespace
\S\ref{sec:spread} refuted inference
 & one inference had been written in \textbf{four} places before being refuted
 & swept repo-wide by equal-length replacement, not patched at the site of
   discovery \\
\addlinespace
\S\ref{sec:fake} a guard that cannot fail
 & a guard passed while its guarded quantity was meaningless, printing a
   plausible fake reading
 & guard strengthened to key on the maximum entry; the incident is registered
   as a live instance of the paper's own thesis \\
\bottomrule
\end{tabular}
\caption{The five self-audit outcomes. None of them changes a physical readout;
all of them change what the ledger is entitled to say.}
\label{tab:selfaudit}
\end{table}

\subsection{Closing a gap is not the same currency as adding evidence}
\label{sec:gapclosure}

The companion paper's ledger had a single uncovered layer, $L1$: it had no
negative control. v16 closed that gap --- coverage is now $7/7$, gap \code{[]},
with 13 controls distributed as 6 repulsion, 1 null-model, 5 two-path and 1
upper-bound.

The natural next sentence is ``therefore $L1$ is promoted''. We did not write it.
The reason is visible only by reading the control's implementation: the
repulsion control evaluates
$\bigl|\langle R_y(\theta)\psi^{\otimes L}\mid g_{\mathrm{ED}}\rangle\bigr|^2$,
where the trial state is prepared by \emph{the same} \code{np.kron} construction
that the primary criterion uses. A control built from the criterion's own
construction cannot falsify that criterion; it can only calibrate the scale on
which the criterion's sensitivity is measured. That property is real and worth
having --- and it is recorded in the control's own note --- but it is not
independent evidence.

The layer therefore remains at its previous maturity, and the ledger records the
promotion as \emph{considered and declined}, with reasons.

\begin{quote}\small
\textbf{Rule.} ``The gap is closed'' and ``the evidence is now independent'' are
different currencies. A coverage counter cannot tell them apart, because both
increment it. Only reading the control's construction can.
\end{quote}

\subsection{The method caught a derivation that was a tautology}
\label{sec:tautology}

One registered item claimed to derive a uniqueness statement about a measure.
Auditing it against the code showed the ``derivation'' to be two ready-made
theorems placed side by side: that the uniform distribution maximises Shannon
entropy, and that the Haar measure is the unique unitary-invariant measure.
Juxtaposition is not derivation, and neither theorem constrains the other.

The audit also found a mismatch in the \emph{object}: the claim's entropy is the
Shannon entropy of a probability vector, whereas the code's state is a pure
state, whose von Neumann entropy is identically zero --- a \emph{minimum}, not a
maximum, of that functional. The two are not the same quantity.

Disposition: registered as a tautology, with the explicit statement that a
tautology \textbf{does not narrow the gap} it was invoked to close. This is the
same disposition the companion paper gave its own reductive results, now applied
reflexively.

\begin{quote}\small
\textbf{Rule.} A ``derivation'' that only re-labels structure already assumed
contributes zero. Audit the derivation's object, not just its conclusion.
\end{quote}

\subsection{The method caught two stale registrations, one of them about the
denominator}
\label{sec:stale}

Two registry entries were found to describe code behaviour that does not exist.

\paragraph{(a) A verdict recorded backwards.} A ledger comment stated that two
checks, $V6$ and $V7$, are both registered with \code{expect_pass=False} and
therefore excluded from the pass rate. Reading the call sites shows that only
$V6$ carries that flag; $V7$ is scored. The registry had been describing a
\emph{verdict} that the code does not implement. This is not cosmetic: had the
comment been taken as authoritative and ``corrected'' in the other direction ---
by adding the flag to $V7$ --- the scored denominator would have dropped from
57 to 56.

\paragraph{(b) A guard count settled by code, not by a log line.} Two documents
and one script banner stated the guard total as 77. The frozen run's report line
(Sec.~\ref{sec:split}) gives $57$ scored $+\,21$ diagnostics $=\,78$, and the
partition function splits a single list --- so $78$ is a consequence of the
current code, whereas $77$ was a consequence of the code as it stood at an
earlier date.

This is the point at which the discipline acquired a rule about \emph{tense}. An
audit found ``77'' in eight places. Four of them are \emph{dated archive
records} --- entries explicitly labelled with an earlier version and date, in
which 77 was the number genuinely measured at the time. Those four were left
\emph{unchanged}. Editing them to 78 would not be a correction; it would be
falsification of the record. The other four were standing assertions about the
current value and were corrected.

\begin{quote}\small
\textbf{Rule.} The same digit may be stale or true depending on tense. A dated
archive assertion is never ``updated''; a present-tense assertion is never left
stale. Before changing a count, broaden the search: the registered list of sites
may itself be miscategorised, and correcting only the listed subset leaves two
mutually contradictory documents.
\end{quote}

\subsection{The method caught a refuted inference in four places}
\label{sec:spread}

An earlier inference --- that a particular measurement \emph{established} that a
certain carrier dimension is not excluded by a no-go argument --- was refuted
when the measurement was recomputed. The refutation was local, but the inference
was not: it had been written into four separate documents, including a runner's
explanatory note and a script's header.

Disposition: the four sites were corrected in one batch, by \emph{equal-length
replacement} --- substituting $N$ lines for $N$ lines --- because one of the
affected files is cited by line number from a smoke test, and inserting a line
would silently break that citation.

\begin{quote}\small
\textbf{Rule.} An inference refuted at one site is not refuted; it is refuted
only when every site carrying it is swept. And when the swept files are cited by
line number, the sweep must be length-preserving, or it will manufacture a
second class of error while fixing the first.
\end{quote}

That second class duly appeared, and is registered in turn: a mid-file insertion
into a document moved the line numbers of citations \emph{into that same
document}, and a section inserted at the \emph{top} of a changelog moved a set
of line-number citations by roughly 150 lines. The content was correct
throughout; only the pointers drifted. The disposition adopted is to convert
fragile line-number citations into verbatim-fragment citations, and to accept
that the already-drifted pointers are registered rather than chased --- since
editing the body to realign them would move other lines, and recurse.

\subsection{A guard that could not fail printed a plausible fake reading}
\label{sec:fake}

This is the paper's centrepiece, and it happened inside the project's own
reference implementation, while it was building a new downstream readout for
$L1$.

\paragraph{The setting.} The readout propagates a trial state with the
imaginary-time operator $e^{-\tau H}$ --- $H$ the periodic transverse-field
Ising Hamiltonian at $L{=}16$ --- and measures the half-chain entropy $S(L/2)$
downstream. The natural time scale is $\tau^\ast = 1/|E_0/L|$; with the
registered $E_0/L = -1.275287$, $\tau^\ast = 0.784137$.

\paragraph{The trap.} $e^{-\tau H}$ is \emph{upward} propagation in imaginary
time: the norm grows as $e^{\tau|E_0|}$. With $|E_0| = 20.4046$ at $L{=}16$, the
growth factor at $\tau{=}20$ is $e^{408}$, and the vector's largest entry is
around $2\times10^{176}$ --- comfortably representable. The implementation's
guard was

\begin{quote}\small
\code{if not np.all(np.isfinite(v)): \# report unavailable}
\end{quote}

\noindent which \emph{passes}, since every entry is finite. Normalisation was
then \code{u = v / np.linalg.norm(v)}. But \code{np.linalg.norm} squares before
it sums; with entries of order $10^{176}$ the square overflows to \code{inf},
the norm is \code{inf}, and \code{v / inf} is the \emph{zero vector}. The
subsequent entropy and overlap are computed from that zero vector and printed:
\[
S(L/2) = 0, \qquad \bigl|\langle g_{\mathrm{gs}} \mid u\rangle\bigr| = 0 .
\]
Both are finite, well-typed, and formatted like every other reading in the
report. Nothing in the output distinguishes this from a genuine result.

\paragraph{The fix.} Guard on the \emph{maximum entry} rather than the norm, and
scale by it before normalising --- a quantity that has not yet been squared, and
so cannot overflow while the entries themselves remain representable. The
overflow bound then becomes a readable fact rather than a crash:
$\tau > 709.78/|E_0| = 34.79$.

\paragraph{Why this belongs in a methodology paper.} The companion paper's
thesis is that a sanity check is consistent with any conclusion, so a wall of
green checks does not distinguish a correct pipeline from a broken one. Here the
check in question --- finiteness --- is precisely a member of the excluded
family, it was written by an author who had stated the exclusion criterion, and
it \emph{passed} while its guarded quantity was meaningless. The output was not
an error message; it was a number.

The episode is registered as three numerical facts, verbatim from the run's own
transcript: at $\tau{=}10$, $\max|v| = 5.063\times10^{87}$; at $\tau{=}20$,
$2.092\times10^{176}$; at $\tau{=}30$, $8.641\times10^{264}$ --- each with the
guarded readings that replaced the fake one. We note that the same episode
supplies a second, smaller instance of Sec.~\ref{sec:stale}'s rule: an author
comment stating the tolerance rationale was itself initially wrong (it quoted a
half-last-digit derived from a four-decimal printed value rather than the
three-decimal value actually printed), and was corrected before the run was
registered.

\begin{quote}\small
\textbf{Rule.} A guard must be keyed on a quantity that \emph{cannot} remain
finite while the guarded value becomes meaningless. ``All entries are finite''
is not such a quantity when the next operation is a sum of squares.
\end{quote}

\subsection{What the five share}
\label{sec:commontheme}

Each of the five is an instance of the same structure: \emph{a check that
incremented correctly while the thing it stood for did not obtain}.

\begin{itemize}[leftmargin=1.4em,itemsep=2pt]
  \item Coverage incremented while independence did not.
  \item A derivation was present while a derivation did not occur.
  \item A registry entry existed while its described behaviour did not.
  \item A correction was applied while the refuted claim survived elsewhere.
  \item A finiteness check passed while the value was meaningless.
\end{itemize}

The composite lesson is that the failure modes are \emph{bookkeeping-shaped},
not arithmetic-shaped. None of the five would be caught by a numerical
tolerance, a longer run, or a larger $L$. All five were caught by asking what
the check's increment \emph{means}.

% ---------------------------------------------------------------------------
\section{Negative controls: which ones can actually fail}
\label{sec:controls}

A control that has never been observed to fail is indistinguishable from a
control that cannot. v16 therefore tested its own controls by scanning them over
declared grids.

\subsection{Four orders of magnitude between a check and its downstream}

For the $L4$ layer, an isometry-destroying perturbation of strength $\varepsilon$
was injected along an isolated path (deliberately bypassing the main entry
point, whose registered readout must not move). Result: a perturbation of
$\varepsilon = 10^{-6}$ already drives the isometry residual above $10^{-12}$,
whereas the downstream central-charge readout does not degrade until
$\varepsilon \approx 3\times10^{-2}$ --- a gap of about four orders of
magnitude. Two further properties are registered alongside, because reporting
only the first would overstate the control: the response is \emph{non-monotone}
(at $\varepsilon = 10^{-3}$ the downstream figure is $0.787\times$ its baseline,
i.e.\ \emph{below} it), and the absolute value of the extracted central charge is
not comparable to production, since the perturbed path runs at a different step
count. It is the \emph{gap} that is the finding, not the endpoint value.

\begin{quote}\small
\textbf{Rule.} A sensitivity margin is a quantitative result and should be
reported as one --- together with any non-monotonicity, which a single ``it
fails'' summary would hide. The margin must be reported on the same page as the
control, or the control will be read as stronger than it is.
\end{quote}

\subsection{Where the discriminating power actually comes from}

For the $L6$ reaction--diffusion layer, the control's two free parameters were
scanned over a $30$-point grid, frozen \emph{before} the scan, with the freeze
line recorded in the log. On that grid the criterion reports ``not emerged'' at
23 of the 30 points, and all 30 are flagged finite --- so the criterion can fail,
and its failure is not an artefact of a numerical breakdown.

Re-partitioning the already-recorded table (no new data, no point selection)
localises the power further: all 23 failing points have a contrast of exactly
zero, whereas only 11 of them also satisfy the activity bound --- the other 12
are uniformity-saturated states. On this grid, discrimination comes from the
contrast criterion \emph{alone}. We state this for this declared grid only.

\subsection{A downstream path that exists, and can fail, and is still not
evidence}
\label{sec:downstream}

The $L1$ layer admits a directly relevant test: propagate the layer's own state
with $e^{-\tau H}$ and ask whether the downstream readout can distinguish it. It
can. At $\tau^\ast$, the correct state gives $S(L/2) = 0.5090065518$ while a
deliberately wrong state gives $0.7150567405$, a relative separation of
$2.88\times10^{-1}$ against a pre-registered threshold of $10^{-3}$. The four
registered readouts of the reference run are reproduced at relative deviations
of $1.6\times10^{-7}$, $4.5\times10^{-5}$, $3.9\times10^{-4}$ and
$4.0\times10^{-7}$, and the entropy readout agrees with the pipeline's own
reference value to $3.3\times10^{-16}$. The acceptance
criterion --- ``it must be possible to fail at least once'' --- is met.

It is still not evidence for the layer, for three reasons, two of which were
written down \emph{before} the measurement:

\begin{enumerate}[leftmargin=1.4em,itemsep=2pt]
  \item The operator acts on a \emph{trial state}; the production readout is the
        \emph{ground state} of $H$. The path is a parallel construction, not the
        production readout, and its agreement therefore does not transfer.
  \item ``Not sharing a construction'' holds only in how the state is written.
        The tensor is assembled by pointwise multiplication and reshaping rather
        than by \code{np.kron} recursion --- verified by counting \code{kron}
        attribute accesses in the syntax tree, giving zero --- but the two are
        \emph{numerically identical}. What is genuinely not shared is the
        readout mechanism, not the state preparation.
  \item The layer's promotion conditions are conjunctive, and this addresses
        only one of them.
\end{enumerate}

\begin{quote}\small
\textbf{Rule.} ``The check can fail'' and ``the check tests the right thing'' are
independent properties. Passing the first while missing the second is the normal
case, not the exceptional one.
\end{quote}

% ---------------------------------------------------------------------------
\section{The one positive template: a criterion that could have chosen otherwise}
\label{sec:positive}

Five of the self-audit outcomes above are negative. It is worth recording that
the project contains exactly one result of the opposite kind, because it
identifies what the negative ones lack.

In the cMERA branch ($L5$), a self-check of the Gaussian representation passes
$9/9$. What distinguishes it is not the pass rate but a structural property of
how it was built: \emph{the criterion selected the convention, rather than the
convention determining the criterion}. The check was fixed first, and the
representation convention was then the thing that had to satisfy it --- so a
different convention would have failed, which is exactly what the other five
results cannot say about themselves.

Two boundaries must be attached, and they are why this section is short.

\paragraph{(i) It is a statement about the criterion.} It says that a convention
was constrained, not that a spacetime emerged.

\paragraph{(ii) What it does not support.} A companion probe found that under the
declared convention the cMERA flow's terminal state is
\emph{indistinguishable} from the bare ultraviolet vacuum, with a relative span
of $2.285\times10^{-15}$ across the objects compared. That is a negative result
about the branch, and it is why $L5$'s maturity is unchanged. A related probe's
original question turned out to have \emph{no observable at all} under the
declared convention --- there were no layers, and hence no per-layer generator
to compare --- so it could not be answered as posed.

\begin{quote}\small
\textbf{The constructive rule.} A result is a result, rather than a tautology,
exactly when the criterion could have come out otherwise. This is the property
whose absence the self-audit outcomes of Sec.~\ref{sec:selfaudit} detect.
\end{quote}

% ---------------------------------------------------------------------------
\section{Boundaries}
\label{sec:boundaries}

We state these explicitly, in the spirit of the ledger itself.

\begin{enumerate}[leftmargin=1.4em,itemsep=2pt]
  \item \textbf{No physical claim.} Nothing here bears on whether an emergent
        spacetime picture \cite{maldacena} is correct. The one anchored
        substrate result --- that the ground-state energy density follows the
        parameter-free finite-size CFT form, and that the central charge
        extracted from a low-energy spectral ratio sits within $0.212\%$ of
        $1/2$ --- is carried over unchanged from the companion paper, where it
        is explicitly noted that the estimate does \emph{not} converge
        monotonically in $L$, and that $L{=}16$ is its worst point.
  \item \textbf{No layer was promoted.} The maturity vector is unchanged from
        the companion paper. In particular, the layer whose coverage gap was
        closed is \emph{not} upgraded, for the reason given in
        Sec.~\ref{sec:gapclosure} --- and that reason was registered as a
        decision rather than discovered as a result.
  \item \textbf{``$100\%$'' means guard execution.} $57$ scored checks passed.
        It is not evidence that the system became better understood.
  \item \textbf{Diagnostics are not failures.} Of the $21$ checks held out of
        the denominator, $16$ do not pass. That is by design: they are checks
        whose guarded proposition is already known to be false, or is not
        threshold-bearing. Excluding them is a convention, and it is stated.
  \item \textbf{Timing is registered, not explained.} The audit records nine
        entries in which readouts agree between runs while wall-clock times do
        not. Those were neither re-measured nor adjusted, and no explanation is
        offered here. A report that invented a mechanism for its own timing
        anomalies would be committing the same category of error this paper
        warns about.
  \item \textbf{Scope of the controls.} Where a control's discriminating power
        is stated, it is stated for the \emph{declared grid} only --- frozen in
        writing before the scan, with the freeze line in the log. No
        extrapolation is claimed.
\end{enumerate}

% ---------------------------------------------------------------------------
\section{Artefacts and reproducibility}
\label{sec:artefacts}

The pipeline is driven by a single runner; all figures and the JSON data
artefact are regenerated from scratch, and no pre-existing data file is read.
The v16 line's incremental probes (five of them) are deliberately kept
\emph{outside} the runner's registry: they are diagnostics, and adding them
would move the denominators reported in Table~\ref{tab:ledger}. Each carries a
version tag and writes its own transcript.

Total elapsed time for the frozen reference run is $1119.7$\,s on an 8-core
machine at $49.9\%$ aggregate CPU utilisation, as printed by the runner itself.

% ---------------------------------------------------------------------------
\section{Conclusion}
\label{sec:conclusion}

The companion paper argued that a verification discipline is defined by whether
its checks can fail. This paper applies that discipline to the artefact
enforcing it, and reports five failures --- of a coverage counter, of a
derivation, of a registry entry, of a sweep, and of a guard --- each of which
had been passing.

The most transferable of the five is the last. A finiteness check, written in a
codebase that explicitly excludes finiteness checks from its ledger, passed on a
vector whose norm had silently overflowed; the program printed $S(L/2) = 0$ and
no warning. There is no tolerance, step count, or system size at which that
becomes visible. It becomes visible only when one asks what the check's
increment means --- which is the whole of the method, and the reason it is worth
applying to one's own work first.

We record the ledger's state rather than a conclusion: 27 metrics, 78 guards, 13
negative controls, $57/57$ scored checks passing, no layer promoted. The value
of the iteration is not in the pass rate. It is in the entries written in the
negative.

% ---------------------------------------------------------------------------
\begin{thebibliography}{9}

\bibitem{v15paper}
J.~Yan,
\newblock A negative-control ledger for multi-stage numerical pipelines
--- demonstrated on Ising/MERA,
\newblock preprint, \code{v15/spiral_v15_preprint.tex}, 2026.

\bibitem{pfeuty}
P.~Pfeuty,
\newblock The one-dimensional Ising model with a transverse field,
\newblock \emph{Ann.\ Phys.}\ \textbf{57}, 79 (1970).

\bibitem{vidal}
G.~Vidal,
\newblock Entanglement renormalization,
\newblock \emph{Phys.\ Rev.\ Lett.}\ \textbf{99}, 220405 (2007).

\bibitem{haegeman}
J.~Haegeman, T.~J.~Osborne, H.~Verschelde, F.~Verstraete,
\newblock Entanglement renormalization for quantum fields in real space,
\newblock \emph{Phys.\ Rev.\ Lett.}\ \textbf{110}, 100402 (2013).

\bibitem{calabrese}
P.~Calabrese and J.~Cardy,
\newblock Entanglement entropy and quantum field theory,
\newblock \emph{J.\ Stat.\ Mech.}\ (2004) P06002.

\bibitem{pearson}
J.~E.~Pearson,
\newblock Complex patterns in a simple system,
\newblock \emph{Science}\ \textbf{261}, 189 (1993).

\bibitem{ollivier}
Y.~Ollivier,
\newblock Ricci curvature of Markov chains on metric spaces,
\newblock \emph{J.\ Funct.\ Anal.}\ \textbf{256}, 810 (2009).

\bibitem{forman}
R.~Forman,
\newblock Bochner's method for cell complexes and combinatorial Ricci curvature,
\newblock \emph{Discrete Comput.\ Geom.}\ \textbf{29}, 323 (2003).

\bibitem{popper}
K.~R.~Popper,
\newblock \emph{The Logic of Scientific Discovery},
\newblock Hutchinson, London, 1959.

\bibitem{maldacena}
J.~Maldacena and L.~Susskind,
\newblock Cool horizons for entangled black holes,
\newblock \emph{Fortsch.\ Phys.}\ \textbf{61}, 781 (2013).

\end{thebibliography}

\end{document}
